\section{Single-bond commuting form from the local structure}

\begin{definition}[Translation-invariant two-site bond]
    \label{def:mpdo_translation_invariant_bond_data}
    \lean{MPOTensor.TranslationInvariantBondData}
    \leanok
    \uses{def:mpdo_commuting_form_data}
    This records one positive semidefinite two-site bond $B\ge0$ whose
    translated copies commute pairwise on the periodic $N$-site chain at
    every chain length $N\ge2$: $[B_{i,i+1},B_{j,j+1}]=0$.
\end{definition}

\begin{definition}[$\eta$-local single-bond commuting-form data]
    \label{def:mpdo_eta_local_structure_data}
    \lean{MPOTensor.EtaLocalStructureData}
    \leanok
    \uses{def:mpdo_translation_invariant_bond_data, def:mpdo_has_commuting_form}
    This defines the one-bond eta-local form constructed in
    \cite[Appendix~C.2]{Cirac2016MPDO_arXiv}:
    a two-site bond $B$ as in
    Definition~\ref{def:mpdo_translation_invariant_bond_data} together with,
    for every chain length $N\ge2$, a constant $c>0$ such that
    \begin{align}
        \rho^{(N)}(M)
        &=c\prod_{i=0}^{N-1}B_{i,i+1}.
        \notag
    \end{align}
    The bond acts on the full two-site space.
\end{definition}

\begin{definition}[Finite-chain bond form]
    \label{def:mpdo_translation_invariant_bond_to_commuting_form_data}
    \lean{MPOTensor.TranslationInvariantBondData.toCommutingFormData}
    \lean{MPOTensor.TranslationInvariantBondData.toCommutingFormData_bond}
    \lean{MPOTensor.TranslationInvariantBondData.toCommutingFormData_hN}
    \lean{MPOTensor.TranslationInvariantBondData.toCommutingFormData_bondAt}
    \leanok
    \uses{def:mpdo_translation_invariant_bond_data, def:mpdo_commuting_form_data}
    A translation-invariant positive two-site bond determines finite-chain
    commuting-form data at every length $N\ge2$.
\end{definition}

\begin{definition}[Common three-site coordinates for overlapping bonds]
    \label{def:mpdo_three_site_overlapping_coordinates}
    \lean{MPOTensor.threeSiteOverlappingEquiv}
    \lean{MPOTensor.pairBondMatrix}
    \lean{MPOTensor.TranslationInvariantBondData.pairBond}
    \lean{MPOTensor.EtaLocalStructureData.pairBond}
    \leanok
    \uses{def:mpdo_translation_invariant_bond_data}
    Write the two-site bond on the ordered pair of physical indices, and
    identify a three-site configuration $(x_0,x_1,x_2)$ with the
    left-associated triple $((x_0,x_1),x_2)$. These are common coordinates
    for the two bonds on sites $(0,1)$ and $(1,2)$.
\end{definition}

\begin{theorem}[Adjacent translates are the two overlapping lifts]
    \label{thm:mpdo_eta_bond_overlapping_lifts}
    \lean{MPOTensor.reindex_embedLocalOperator_zero_eq_leftOverlappingLift}
    \lean{MPOTensor.reindex_embedLocalOperator_one_eq_rightOverlappingLift}
    \lean{MPOTensor.TranslationInvariantBondData.pairBond_isHermitian}
    \lean{MPOTensor.TranslationInvariantBondData.%
      reindex_bondAt_zero_eq_leftOverlappingLift}
    \lean{MPOTensor.TranslationInvariantBondData.%
      reindex_bondAt_one_eq_rightOverlappingLift}
    \lean{MPOTensor.TranslationInvariantBondData.%
      overlappingLifts_pairBond_comm}
    \lean{MPOTensor.EtaLocalStructureData.pairBond_isHermitian}
    \lean{MPOTensor.EtaLocalStructureData.reindex_bondAt_zero_eq_leftOverlappingLift}
    \lean{MPOTensor.EtaLocalStructureData.reindex_bondAt_one_eq_rightOverlappingLift}
    \lean{MPOTensor.EtaLocalStructureData.overlappingLifts_pairBond_comm}
    \leanok
    \uses{def:mpdo_three_site_overlapping_coordinates,
        def:mpdo_translation_invariant_bond_data}
    Let $B$ be a translation-invariant positive two-site bond, written as an
    operator $\widehat B$ on an ordered pair of physical spaces.
    Under the same three-site identification for both translates,
    \begin{align}
        B_{0,1}
        &=\widehat B\otimes\Id,
        \notag\\
        B_{1,2}
        &=\Id\otimes\widehat B.
        \notag
    \end{align}
    Thus these two operators are respectively the left and right overlapping
    lifts of $\widehat B$. Moreover, $\widehat B$ is Hermitian and the two
    lifts commute. These are exactly the local hypotheses needed for the
    Bravyi--Vyalyi decomposition invoked at
    \cite[Appendix~C.2, lines~1599--1605]{Cirac2016MPDO_arXiv}.
\end{theorem}

\begin{proof}\leanok
    Let $\Phi$ be the matrix reindexing induced by the common three-site
    identification, and write $B_{0,1}=\Phi(B_0)$ and
    $B_{1,2}=\Phi(B_1)$ for the two reindexed translates. For three-site
    configurations $x$ and $y$, the two outside-window conditions are
    \begin{align}
        (\forall r\notin\{0,1\},\ x_r=y_r)
        &\Longleftrightarrow x_2=y_2,
        \notag\\
        (\forall r\notin\{1,2\},\ x_r=y_r)
        &\Longleftrightarrow x_0=y_0.
        \notag
    \end{align}
    Hence evaluation at a pair of left-associated triples gives
    \begin{align}
        (B_{0,1})_{(i,j,k),(i',j',k')}
        &=\widehat B_{(i,j),(i',j')}\delta_{k,k'}
          =(\widehat B\otimes\Id)_{(i,j,k),(i',j',k')},
        \notag\\
        (B_{1,2})_{(i,j,k),(i',j',k')}
        &=\delta_{i,i'}\widehat B_{(j,k),(j',k')}
          =(\Id\otimes\widehat B)_{(i,j,k),(i',j',k')}.
        \notag
    \end{align}
    Positivity of $B$ gives $\widehat B^*=\widehat B$. Finally, matrix
    reindexing is multiplicative, so $\Phi(XY)=\Phi(X)\Phi(Y)$, and
    \begin{align}
        B_0B_1=B_1B_0
        &\Longrightarrow
        \Phi(B_0)\Phi(B_1)=\Phi(B_1)\Phi(B_0).
        \notag
    \end{align}
    Substituting the two entrywise identities proves that the left and right
    overlapping lifts commute.
\end{proof}

\begin{theorem}[$\eta$-local structure carries single-bond commuting form]
    \label{thm:mpdo_eta_local_structure_has_commuting_form}
    \lean{MPOTensor.EtaLocalStructureData.hasCommutingForm}
    \lean{MPOTensor.EtaLocalStructureData.formAt_realizes}
    \lean{MPOTensor.EtaLocalStructureData.formAt_bondAt_comm}
    \lean{MPOTensor.EtaLocalStructureData.exists_positive_scalar_mpo_eq_product}
    \lean{MPOTensor.EtaLocalStructureData.positive_commuting_product_form}
    \lean{MPOTensor.hasCommutingForm_of_etaLocalStructure}
    \leanok
    \uses{def:mpdo_eta_local_structure_data, def:mpdo_has_commuting_form}
    The $\eta$-local structure itself determines a single-bond
    commuting-form witness for every chain length $N\ge2$: writing $B$ for
    its bond, the translated copies commute,
    $[B_{i,i+1},B_{j,j+1}]=0$, and there exists $c>0$ with
    \begin{align}
        \rho^{(N)}(M)
        &=c\prod_{i=0}^{N-1}B_{i,i+1}.
        \notag
    \end{align}
    This is the one-bond product conclusion used in
    \cite[Proposition~C.8]{Cirac2016MPDO_arXiv}.
    % Gap documented in docs/paper-gaps/cpgsv17_mpdo_sal_zcl_eta_local_structure.tex.
\end{theorem}

\begin{proof}\leanok
    At length $N$, apply the commuting two-site bond and realization identity
    provided by the $\eta$-local structure.
\end{proof}

\begin{theorem}[$\eta$-local structure gives a single-bond product]
    \label{thm:mpdo_eta_local_structure_has_commuting_bond_product}
    \lean{MPOTensor.EtaLocalStructureData.hasCommutingBondProductAt}
    \lean{MPOTensor.EtaLocalStructureData.hasCommutingBondProduct}
    \lean{MPOTensor.hasCommutingBondProduct_of_etaLocalStructure}
    \leanok
    \uses{def:mpdo_eta_local_structure_data,
        def:mpdo_has_commuting_bond_product}
    The same $\eta$-local structure gives a single-bond product witness
    on every finite periodic chain.
\end{theorem}

\begin{proof}\leanok
    Evaluate the local structure at each chain length $N\ge2$.
\end{proof}

\begin{theorem}[Single-bond product with doubled-index idempotence]
    \label{thm:mpdo_eta_local_structure_has_commuting_bond_product_zcl}
    \lean{MPOTensor.EtaLocalStructureData.hasCommutingBondProductWithZCL}
    \lean{MPOTensor.hasCommutingBondProductWithZCL_of_etaLocalStructure}
    \leanok
    \uses{thm:mpdo_eta_local_structure_has_commuting_form,
        thm:mpdo_has_commuting_bond_product_zcl_iff_has_commuting_form_and_is_zcl}
    An $\eta$-local structure together with the doubled-index transfer
    condition gives the corresponding single-bond branch.
    % Scope restriction (single-bond presentation): see
    % docs/paper-gaps/cpsv16_gsnnch_sector_decomposition.tex.
\end{theorem}

\begin{proof}\leanok
    Use the commuting-form witness carried by the $\eta$-local structure and
    combine it with the doubled-index transfer hypothesis.
\end{proof}

\begin{theorem}[A commuting bond admits unitary block actions]
    \label{thm:mpdo_eta_local_pair_bond_unitary_block_actions}
    \lean{MPOTensor.TranslationInvariantBondData.%
      exists_unitary_blockActions_of_pairBond}
    \lean{MPOTensor.EtaLocalStructureData.exists_unitary_blockActions_of_pairBond}
    \leanok
    \uses{def:mpdo_three_site_overlapping_coordinates,
        thm:mpdo_eta_bond_overlapping_lifts,
        thm:commuting_overlapping_unitary_block_actions}
    Let $\widehat B$ be a pair-indexed positive two-site bond whose periodic
    translates commute. There are positive dimensions
    $d_{l,q}=\dim H_{q,l}$ and $d_{r,q}=\dim H_{q,r}$, a unitary decomposition
    $H\cong\bigoplus_{q=0}^{K-1}H_{q,l}\otimes H_{q,r}$,
    and Hermitian operators $R_q$ on $H\otimes H_{q,l}$ and $S_q$ on
    $H_{q,r}\otimes H$ such that
    \begin{align}
        (\Id_H\otimes U^*)\widehat B(\Id_H\otimes U)
        &=\bigoplus_{q=0}^{K-1}
          (R_q\otimes\Id_{H_{q,r}}),
        \notag\\
        (U^*\otimes\Id_H)\widehat B(U\otimes\Id_H)
        &=\bigoplus_{q=0}^{K-1}
          (\Id_{H_{q,l}}\otimes S_q).
        \notag
    \end{align}
    The basis bijection for each sector enumerates the numerical indices in
    right--left order,
    \begin{align}
        \bigsqcup_q\bigl(\{0,\ldots,d_{r,q}-1\}
            \times\{0,\ldots,d_{l,q}-1\}\bigr)
        &\cong\{0,\ldots,d-1\}.
        \notag
    \end{align}
    Writing $e(q,(r,s))$ for this basis bijection, the left and right spatial
    coordinates are respectively
    \begin{align}
        (q,((i,s),r))
        &\longmapsto(i,e(q,(r,s))),
        \notag\\
        (q,(s,(r,k)))
        &\longmapsto(e(q,(r,s)),k).
        \notag
    \end{align}
    Hence the tensor factors in the displayed block actions occur in
    left--right order.
\end{theorem}

\begin{proof}\leanok
    \uses{thm:mpdo_eta_bond_overlapping_lifts,
        thm:commuting_overlapping_unitary_block_actions}
    The pair-indexed bond is Hermitian by positivity, and its two overlapping
    lifts commute by the common three-site identification. Apply
    Theorem~\ref{thm:commuting_overlapping_unitary_block_actions}.
\end{proof}

\begin{theorem}[Positive neighboring operators from the commuting bond]
    \label{thm:mpdo_commuting_bond_positive_eta_decomposition}
    \lean{MPOTensor.TranslationInvariantBondData.%
      exists_positive_eta_pairBond_decomposition}
    \lean{MPOTensor.EtaLocalStructureData.%
      exists_positive_eta_pairBond_decomposition}
    \leanok
    \uses{thm:mpdo_eta_local_pair_bond_unitary_block_actions}
    Let $B$ be a pair-indexed positive bond whose periodic translates commute.
    There are a unitary coordinate decomposition
    $H\cong\bigoplus_q H_{q,r}\otimes H_{q,l}$
    and positive semidefinite operators
    $\eta_{q,h}$ on $H_{q,r}\otimes H_{h,l}$. After regrouping the
    two-site coordinates, the $(q,h)$-sector is
    $H_{q,l}\otimes(H_{q,r}\otimes H_{h,l})\otimes H_{h,r}$,
    and in these regrouped coordinates
    \begin{align}
        (U^*\otimes U^*)B(U\otimes U)
        &=\bigoplus_{q,h}
          \Id_{H_{q,l}}\otimes\eta_{q,h}\otimes\Id_{H_{h,r}}.
        \notag
    \end{align}
    The unitary and the neighboring operators depend only on the fixed bond,
    and hence are independent of a chain length or bond position.
\end{theorem}

\begin{proof}\leanok
    \uses{thm:mpdo_eta_local_pair_bond_unitary_block_actions}
    The same unitary block decomposition at both ends of $B$ gives identity
    action on $H_{q,l}$ through one block identity and on $H_{h,r}$ through
    the other. For fixed basis indices
    $\ell_q^0\in H_{q,l}$ and $r_h^0\in H_{h,r}$, let
    $B'=(U^*\otimes U^*)B(U\otimes U)$. The neighboring operator has entries
    \begin{align}
        (\eta_{q,h})_{(r,\ell),(r',\ell')}
        &=
        B'_{(q,r,\ell_q^0;h,r_h^0,\ell),
          (q,r',\ell_q^0;h,r_h^0,\ell')}.
        \notag
    \end{align}
    Thus $\eta_{q,h}$ is a principal compression of the positive operator
    $B'$, and is therefore positive semidefinite.
\end{proof}

\begin{theorem}[Cyclic product of a locally decomposed bond]
    \label{thm:mpdo_eta_pair_bond_cyclic_product_transport}
    \lean{MPOTensor.%
      reindex_product_embedLocalOperator_of_etaPair_decomposition}
    \leanok
    Suppose that, in one fixed coordinate decomposition
    $H\cong\bigoplus_q H_{q,r}\otimes H_{q,l}$,
    after regrouping the two-site coordinates according to this decomposition,
    a two-site operator has the form
    \begin{align}
        B
        &=\bigoplus_{q,h}
          \Id_{H_{q,l}}\otimes\eta_{q,h}\otimes\Id_{H_{h,r}}.
        \notag
    \end{align}
    For every $N\ge2$, regroup the cyclic chain as
    \begin{align}
        H^{\otimes N}
        &\cong
        \bigoplus_{k:\Z/N\Z\to\{0,\ldots,K-1\}}
        \bigotimes_{n\in\Z/N\Z}
          \left(H_{k_n,r}\otimes H_{k_{n+1},l}\right).
        \notag
    \end{align}
    Then, in these coordinates and with $k_N=k_0$,
    \begin{align}
        \prod_{n=0}^{N-1}B_{n,n+1}
        &=
        \bigoplus_k\bigotimes_{n=0}^{N-1}\eta_{k_n,k_{n+1}}.
        \notag
    \end{align}
    At $N=2$, the two translated windows have the opposite orders $(0,1)$
    and $(1,0)$ and give the two oriented factors
    $\eta_{k_0,k_1}$ and $\eta_{k_1,k_0}$.
\end{theorem}

\begin{proof}\leanok
    In a fixed sector configuration, the cyclic regrouping sends the factors
    $(r_n,l_n)$ to the edge factors $(r_n,l_{n+1})$. The translate beginning
    at $n$ acts by $\eta_{k_n,k_{n+1}}$ on the $n$th edge factor and by the
    identity on every other edge factor. Different sector configurations do
    not mix. Multiplying the translated operators therefore gives the stated
    direct sum of tensor products.
\end{proof}

% **Local fix (proportionality scalar):** The converse at source lines
% 1603--1605 invokes the coefficient-free equation sigmaNK2 after assuming
% only proportionality to the commuting-bond product.  The statement below
% retains that scalar.  See docs/paper-gaps/cpsv16_commuting_form_to_sal.tex.
\begin{theorem}[Scaled finite-chain cyclic neighboring-operator form]
    \label{thm:mpdo_eta_local_sigma_nk2}
    \lean{MPOTensor.sitewisePhysicalMatrix_conjTranspose}
    \lean{MPOTensor.reindex_sitewisePhysicalMatrix_two}
    \lean{MPOTensor.sitewisePhysicalMatrix_one}
    \lean{MPOTensor.sitewisePhysicalMatrix_mul_conjTranspose}
    \lean{MPOTensor.singleKrausMap_bondProduct_of_unitary}
    \lean{MPOTensor.EtaLocalStructureData.%
      exists_positive_cyclic_eta_block_decomposition}
    \leanok
    \uses{def:mpdo_eta_local_structure_data,
        thm:mpdo_commuting_bond_positive_eta_decomposition,
        thm:mpdo_eta_pair_bond_cyclic_product_transport}
    Let $\sigma^{(N)}(\mathcal K)$ be realized, for every $N\ge2$, by
    translates of one positive commuting bond. There are one unitary $U$,
    positive sector dimensions, and positive semidefinite operators
    $\eta_{q,h}$, all independent of $N$, such that for every $N\ge2$ there
    is a constant $c_N>0$ satisfying, with $k_N=k_0$,
    \begin{align}
        E_N^*(U^*)^{\otimes N}\sigma^{(N)}(\mathcal K)
          U^{\otimes N}E_N
        &=
        c_N\bigoplus_{k_0,\ldots,k_{N-1}}
          \bigotimes_{n=0}^{N-1}\eta_{k_n,k_{n+1}}.
        \notag
    \end{align}
    Here $E_N$ is the cyclic regrouping from the edge factors
    $H_{k_n,r}\otimes H_{k_{n+1},l}$ to the site factors
    $H_{k_n,r}\otimes H_{k_n,l}$. The factor $c_N$ is retained from the
    proportionality relation in
    \cite[Appendix~C.2, lines~1571--1576]{Cirac2016MPDO_arXiv}. The passage at
    \cite[Appendix~C.2, lines~1603--1605]{Cirac2016MPDO_arXiv} suppresses this
    factor when invoking the coefficient-free equation \emph{sigmaNK2}; the
    coefficient-free assertion is not part of the present theorem.
\end{theorem}

\begin{proof}\leanok
    \uses{thm:mpdo_commuting_bond_positive_eta_decomposition,
        thm:mpdo_eta_pair_bond_cyclic_product_transport}
    Choose the unitary and the neighboring operators from the decomposition
    of the fixed bond. Conjugation by $(U^*)^{\otimes N}$ carries the product
    of its translates to the product of the conjugated two-site bonds. The
    cyclic regrouping theorem changes this product into the displayed direct
    sum. Applying both transformations to
    $\sigma^{(N)}(\mathcal K)=c_N\prod_n B_{n,n+1}$ preserves the same
    positive scalar $c_N$.
\end{proof}

\begin{remark}[An arbitrary proportional presentation need not have unit scalar]
    The proportional commuting-bond form alone does not permit one common
    local normalization. Indeed, take physical dimension $d=1$ and bond
    dimension $D=2$, with sole matrix-product letter $M^{00}=I_2$, and take
    the two-site bond to be the scalar $B=1$. Then
    \begin{align}
        \sigma^{(N)}
        &=\tr(I_2^N)=2
          =c_N\prod_{n=0}^{N-1}B_{n,n+1},
        &c_N&=2.
        \notag
    \end{align}
    Any positive neighboring decomposition of the one-dimensional physical
    space has one sector and one scalar $a\ge0$. A coefficient-free identity
    at lengths two and four would give $a^2=2$ and $a^4=2$, a contradiction.
    This example is not a normal injective block. More generally, the source
    does not claim that every proportional commuting-bond presentation can be
    made coefficient-free by one local rescaling. In the SAL-to-commuting-form
    direction, the proof chooses the neighboring operators from the source
    tensor and obtains scalar one directly. This distinguishes the
    proportionality at
    \cite[Appendix~C.2, lines~1571--1576]{Cirac2016MPDO_arXiv} from the
    coefficient-free invocation at
    \cite[Appendix~C.2, lines~1603--1605]{Cirac2016MPDO_arXiv}. The
    source-selected normal representative is
    Corollary~\ref{cor:mpdo_normal_fixed_product_tensor_of_is_sal}.
\end{remark}

\begin{definition}[Cyclic edge-weight tensor]
    \label{def:mpdo_cyclic_edge_weight_tensor}
    \lean{MPOTensor.cyclicEdgeWeightTensor}
    \lean{MPOTensor.cyclicEdgeWeightTensor_apply_eq}
    \lean{MPOTensor.etaCyclicEdgeWeight}
    \leanok
    A scalar weight $w(i,j;i',j')$ on two consecutive ket--bra pairs
    determines a translation-invariant matrix-product tensor with virtual
    dimension $d^2$. Its incoming virtual pair is constrained to equal the
    current physical ket--bra pair, while its outgoing pair records the
    following physical pair.
\end{definition}

\begin{theorem}[Closed cyclic edge-weight identity]
    \label{thm:mpdo_cyclic_edge_weight_tensor}
    \lean{MPOTensor.mpo_cyclicEdgeWeightTensor}
    \leanok
    \uses{def:mpdo_cyclic_edge_weight_tensor}
    For every positive length $N$, and with
    $\sigma_N=\sigma_0$ and $\tau_N=\tau_0$, the closed matrix-product
    operator of the cyclic edge-weight tensor is
    \begin{align}
        \rho^{(N)}_{\sigma,\tau}
        &=
        \prod_{n=0}^{N-1}
          w(\sigma_n,\tau_n;\sigma_{n+1},\tau_{n+1}).
        \notag
    \end{align}
    In the virtual trace there is precisely one potentially nonzero cyclic
    virtual configuration.
\end{theorem}

\begin{definition}[Fixed tensor for a bond-product family]
    \label{def:mpdo_fixed_bond_product_tensor}
    \lean{MPOTensor.TranslationInvariantBondData.FixedProductTensorData}
    \lean{MPOTensor.TranslationInvariantBondData.CyclicEdgeWeightForm}
    \lean{MPOTensor.TranslationInvariantBondData.CyclicEdgeWeightForm.%
      toFixedProductTensorData}
    \leanok
    \uses{def:mpdo_translation_invariant_bond_data,
        thm:mpdo_cyclic_edge_weight_tensor}
    A fixed tensor for a bond-product family consists of a positive virtual
    dimension and one matrix-product tensor $C$, independent of $N$, such
    that, for every $N\ge2$,
    \begin{align}
        \rho^{(N)}(C)
        &=\prod_{n=0}^{N-1}B_{n,n+1}.
        \notag
    \end{align}
    If $d>0$, a cyclic scalar edge-weight expression for the entries of the
    product gives such a tensor with positive virtual dimension $d^2$.
\end{definition}

\begin{theorem}[Fixed tensor representation of a commuting bond]
    \label{thm:mpdo_eta_fixed_bond_product_tensor}
    \lean{MPOTensor.%
      reindex_mpo_cyclicEdgeWeightTensor_etaCyclicEdgeWeight}
    \lean{MPOTensor.TranslationInvariantBondData.%
      nonempty_fixedProductTensorData}
    \lean{MPOTensor.TranslationInvariantBondData.fixedProductTensorData}
    \lean{MPOTensor.EtaLocalStructureData.nonempty_fixedProductTensorData}
    \lean{MPOTensor.EtaLocalStructureData.fixedProductTensorData}
    \leanok
    \uses{def:mpdo_translation_invariant_bond_data,
        def:mpdo_eta_local_structure_data,
        def:mpdo_fixed_bond_product_tensor,
        thm:mpdo_commuting_bond_positive_eta_decomposition}
    Let $B\ge0$ be a two-site bond whose periodic translates commute. There
    is a positive integer $R$ and a matrix-product tensor $C$ of virtual
    dimension $R$, both independent of the chain length, such that for every
    $N\ge2$,
    \begin{align}
        \rho^{(N)}(C)
        &=\prod_{n=0}^{N-1}B_{n,n+1}.
        \notag
    \end{align}
    The construction gives a positive $R$ independent of $N$. No normality,
    zero-correlation-length, or scalar normalization hypothesis is required.
\end{theorem}

\begin{proof}\leanok
    \uses{thm:mpdo_commuting_bond_positive_eta_decomposition}
    Choose Beigi's chain-independent one-site unitary and the positive
    neighboring operators $\eta_{q,h}$. For a common matrix-unit index set
    $\mathcal A$, write
    \begin{align}
        \eta_{q,h}((r,l),(r',l'))
        &=
        \sum_{a\in\mathcal A}
        R_a(q)_{r,r'}L_a(h)_{l,l'}.
        \notag
    \end{align}
    Thus the contraction of the right factor at one site with the left factor
    at the next satisfies
    $\sum_{a\in\mathcal A}R_a(q)\otimes L_a(h)=\eta_{q,h}$.
    If $V$ is the resulting unitary change to the two-site sector
    coordinates, this contraction gives
    $\widetilde B=VBV^\dagger$ and
    $B_{\mathrm{phys}}=V^\dagger\widetilde B V=B$.
    For a nonzero physical space, $\mathcal A$ is nonempty and gives a positive
    virtual dimension. For the zero-dimensional physical space, take one
    virtual state and no physical sectors. The cyclic contraction formula
    then gives the product of the translates of $B$ at every length $N\ge2$.
\end{proof}

\begin{theorem}[Positive physical sectors of the selected fixed tensor]
    \label{thm:mpdo_fixed_bond_positive_physical_sector_representative}
    \lean{MPOTensor.TranslationInvariantBondData.%
      exists_positive_physicalSectorFactorization_fixedProductTensorData}
    \lean{MPOTensor.TranslationInvariantBondData.%
      PositivePhysicalSectorFixedProductTensorData}
    \lean{MPOTensor.TranslationInvariantBondData.%
      nonempty_positivePhysicalSectorFixedProductTensorData}
    \lean{MPOTensor.TranslationInvariantBondData.%
      positivePhysicalSectorFixedProductTensorData}
    \lean{MPOTensor.TranslationInvariantBondData.%
      fixedProductTensorDataPhysicalSectorFactorization}
    \lean{MPOTensor.TranslationInvariantBondData.%
      fixedProductTensorDataPhysicalSectorFactorization_physicalBond_eq}
    \lean{MPOTensor.TranslationInvariantBondData.%
      fixedProductTensorDataPhysicalSectorFactorization_neighboring_pos}
    \lean{MPOTensor.PhysicalSectorFactorization.etaSectorFinEquiv}
    \lean{MPOTensor.PhysicalSectorFactorization.%
      sectorCoordinateBond_etaPair_decomposition}
    \leanok
    \uses{thm:mpdo_eta_fixed_bond_product_tensor}
    Let $C$ be the selected fixed tensor of the preceding theorem. There are
    physical sectors
    $\C^d\cong\bigoplus_{q=0}^{K-1}H_q^L\otimes H_q^R$,
    a unitary $U$, and matrices $L_\beta(q)$ and $R_\alpha(q)$ such that,
    writing $C_{\beta,\alpha}$ for the physical matrix at fixed virtual
    indices,
    \begin{align}
        U C_{\beta,\alpha}U^\dagger
        &=
        \bigoplus_{q=0}^{K-1}
          L_\beta(q)\otimes R_\alpha(q).
        \notag
    \end{align}
    The neighboring operators
    $\eta_{q,h}=\sum_a R_a(q)\otimes L_a(h)$
    are positive semidefinite. If $U_2$ denotes the induced two-site change
    of coordinates, then
    \begin{align}
        B
        &=
        U_2^\dagger
        \left[\bigoplus_{q,h}
          \Id_{H_q^L\otimes H_h^R}\otimes\eta_{q,h}\right]
        U_2.
        \notag
    \end{align}
    Thus the physical-sector factorization belongs to the same tensor $C$
    whose closed operators are the fixed products of $B$.
\end{theorem}

\begin{proof}\leanok
    Take the physical-sector factorization retained with $C$. Its defining
    one-site identity is
    \begin{align}
        U C_{\beta,\alpha}U^\dagger
        &=
        \bigoplus_q L_\beta(q)\otimes R_\alpha(q).
        \notag
    \end{align}
    Contracting the adjacent right and left factors gives
    $\eta_{q,h}=\sum_a R_a(q)\otimes L_a(h)\succeq0$,
    where positivity is the retained neighboring-operator property. Finally,
    the retained bond identity is
    \begin{align}
        B
        &=
        U_2^\dagger
        \left[\bigoplus_{q,h}
          \Id_{H_q^L\otimes H_h^R}\otimes\eta_{q,h}\right]
        U_2.
        \notag
    \end{align}
\end{proof}

\begin{remark}[The scalar edge formula depends on the one-site coordinates]
    A cyclic scalar edge formula need not hold in the original physical basis.
    For $d=2$, let $X$ be the Pauli flip and set $B=\Id+X\otimes X$. This bond
    is positive semidefinite and its periodic translates commute. For every
    $N\ge2$,
    \begin{align}
        \bra{0^N}
        \left(\prod_n B_{n,n+1}\right)
        \ket{0^N}
        &=2.
        \notag
    \end{align}
    At $N=2$ the periodic product is $B^2=2B$; for $N\ge3$, the empty edge
    set and the full cycle are the two terms with zero boundary. An
    original-coordinate scalar edge weight would make this entry equal to
    $a^N$. The cases $N=2$ and $N=3$ would give $a^2=2$ and $a^3=2$, a
    contradiction. Equation \emph{sigmaNK2} in
    \cite[Appendix~C.2, lines~1581--1589]{Cirac2016MPDO_arXiv} is stated only
    after the one-site unitary coordinate change, as used in the theorem.
\end{remark}

\begin{theorem}[Eventual proportionality from a fixed product tensor]
    \label{thm:mpdo_eta_fixed_product_tensor_eventual_proportionality}
    \lean{MPOTensor.EtaLocalStructureData.%
      eventuallyNonzeroProportionalMPV₂_fixedProductTensor}
    \leanok
    \uses{def:mpdo_eta_local_structure_data,
        def:mpdo_fixed_bond_product_tensor,
        thm:mpdo_eta_fixed_bond_product_tensor}
    Suppose one fixed tensor $C$ generates the products of the bond carried
    by an eta-local structure. Then the doubled-index tensors of $M$ and
    $C$ are eventually nonzero proportional. The proof uses only lengths
    $N\ge2$ and therefore assumes no relation at length one.
\end{theorem}

\begin{proof}\leanok
    At every $N\ge2$, substitute the exact tensor representation into the
    positive realization
    $\rho^{(N)}(M)=c_N\prod_n B_{n,n+1}$, where $c_N>0$, and read ket--bra
    pairs as doubled physical indices.
\end{proof}

\begin{theorem}[A normal rescaling gives a geometric phase]
    \label{thm:mpdo_eta_normal_rescaling_geometric_phase}
    \lean{MPOTensor.EtaLocalStructureData.%
      exists_unit_phase_power_mpo_eq_product_of_normal_smul}
    \leanok
    \uses{def:mpdo_eta_local_structure_data,
        def:mpdo_fixed_bond_product_tensor, def:nt_cpsv}
    Let $A$ be the normal doubled-index source tensor of positive bond
    dimension $D$, and let $C$ be a fixed tensor for its commuting-bond
    product. Suppose that $sC$ is normal for some $s>0$. Then there is a
    complex number $\zeta$, with $|\zeta|=1$, such that, for every $N\ge2$,
    \begin{align}
        \rho^{(N)}(A)
        &=
        (\zeta s)^N\prod_{n=0}^{N-1}B_{n,n+1}.
        \notag
    \end{align}
    This is a conditional reduction: it does not assert that $C$, or any
    rescaling of $C$, is normal.
\end{theorem}

\begin{proof}\leanok
    \uses{thm:mpdo_eta_fixed_product_tensor_eventual_proportionality,
        lem:mpv_smul, thm:normal_proportional_mpv_geometric_phase,
        lem:mpdo_three_first_site_contraction_coefficients}
    Scaling $C$ by $s$ gives
    $V^{(N)}(sC)_\sigma=s^N V^{(N)}(C)_\sigma$.
    Thus, for all sufficiently large $N$, the proportionality
    $V^{(N)}(A)_\sigma=c_NV^{(N)}(C)_\sigma$ becomes
    \begin{align}
        V^{(N)}(A)_\sigma
        &=
        \frac{c_N}{s^N}V^{(N)}(sC)_\sigma,
        &\frac{c_N}{s^N}&\ne0.
        \notag
    \end{align}
    The normal-tensor overlap dichotomy gives one $\zeta\in\C$, with
    $|\zeta|=1$, such that, at every positive length,
    \begin{align}
        V^{(N)}(A)_\sigma
        &=
        \zeta^N V^{(N)}(sC)_\sigma
        =(\zeta s)^N V^{(N)}(C)_\sigma.
        \notag
    \end{align}
\end{proof}

\begin{theorem}[A normal rescaling forces the positive geometric law]
    \label{thm:mpdo_eta_normal_rescaling_positive_geometric_law}
    \lean{MPOTensor.EtaLocalStructureData.%
      mpo_eq_positive_power_product_of_normal_smul}
    \leanok
    \uses{def:mpdo_eta_local_structure_data,
        def:mpdo_fixed_bond_product_tensor, def:nt_cpsv}
    Let $A$ be the normal doubled-index source tensor of positive bond
    dimension $D$, let $C$ be a fixed tensor for its commuting-bond product,
    and suppose that $sC$ is normal for some $s>0$. Then the unit phase is
    trivial. For every $N\ge2$,
    \begin{align}
        \rho^{(N)}(A)
        &=
        s^N\prod_{n=0}^{N-1}B_{n,n+1}.
        \notag
    \end{align}
    Thus a positive normal rescaling of the fixed product tensor already
    supplies the positive geometric realization law.
\end{theorem}

\begin{proof}\leanok
    \uses{thm:mpdo_eta_normal_rescaling_geometric_phase,
        thm:cpsv_normal_mpv_phase_gauge}
    Normality gives
    $\lim_{N\to\infty}\braket{V^{(N)}(A)}{V^{(N)}(A)}=1$.
    Hence at two sufficiently large consecutive lengths, say $N$ and $N+1$,
    both the source vector and the fixed bond product are nonzero. Positivity
    of the realization gives positive numbers $c_N,c_{N+1}$ satisfying
    $(\zeta s)^N=c_N$ and $(\zeta s)^{N+1}=c_{N+1}$. Consequently
    $\zeta s=c_{N+1}/c_N>0$. Since $|\zeta|=1$ and $s>0$, one has
    $\zeta=1$.
\end{proof}
